\documentclass[11pt]{article}

    \usepackage[breakable]{tcolorbox}
    \usepackage{parskip} % Stop auto-indenting (to mimic markdown behaviour)
    

    % Basic figure setup, for now with no caption control since it's done
    % automatically by Pandoc (which extracts ![](path) syntax from Markdown).
    \usepackage{graphicx}
    % Maintain compatibility with old templates. Remove in nbconvert 6.0
    \let\Oldincludegraphics\includegraphics
    % Ensure that by default, figures have no caption (until we provide a
    % proper Figure object with a Caption API and a way to capture that
    % in the conversion process - todo).
    \usepackage{caption}
    \DeclareCaptionFormat{nocaption}{}
    \captionsetup{format=nocaption,aboveskip=0pt,belowskip=0pt}

    \usepackage{float}
    \floatplacement{figure}{H} % forces figures to be placed at the correct location
    \usepackage{xcolor} % Allow colors to be defined
    \usepackage{enumerate} % Needed for markdown enumerations to work
    \usepackage{geometry} % Used to adjust the document margins
    \usepackage{amsmath} % Equations
    \usepackage{amssymb} % Equations
    \usepackage{textcomp} % defines textquotesingle
    % Hack from http://tex.stackexchange.com/a/47451/13684:
    \AtBeginDocument{%
        \def\PYZsq{\textquotesingle}% Upright quotes in Pygmentized code
    }
    \usepackage{upquote} % Upright quotes for verbatim code
    \usepackage{eurosym} % defines \euro

    \usepackage{iftex}
    \ifPDFTeX
        \usepackage[T1]{fontenc}
        \IfFileExists{alphabeta.sty}{
              \usepackage{alphabeta}
          }{
              \usepackage[mathletters]{ucs}
              \usepackage[utf8x]{inputenc}
          }
    \else
        \usepackage{fontspec}
        \usepackage{unicode-math}
    \fi

    \usepackage{fancyvrb} % verbatim replacement that allows latex
    \usepackage{grffile} % extends the file name processing of package graphics
                         % to support a larger range
    \makeatletter % fix for old versions of grffile with XeLaTeX
    \@ifpackagelater{grffile}{2019/11/01}
    {
      % Do nothing on new versions
    }
    {
      \def\Gread@@xetex#1{%
        \IfFileExists{"\Gin@base".bb}%
        {\Gread@eps{\Gin@base.bb}}%
        {\Gread@@xetex@aux#1}%
      }
    }
    \makeatother
    \usepackage[Export]{adjustbox} % Used to constrain images to a maximum size
    \adjustboxset{max size={0.9\linewidth}{0.9\paperheight}}

    % The hyperref package gives us a pdf with properly built
    % internal navigation ('pdf bookmarks' for the table of contents,
    % internal cross-reference links, web links for URLs, etc.)
    \usepackage{hyperref}
    % The default LaTeX title has an obnoxious amount of whitespace. By default,
    % titling removes some of it. It also provides customization options.
    \usepackage{titling}
    \usepackage{longtable} % longtable support required by pandoc >1.10
    \usepackage{booktabs}  % table support for pandoc > 1.12.2
    \usepackage{array}     % table support for pandoc >= 2.11.3
    \usepackage{calc}      % table minipage width calculation for pandoc >= 2.11.1
    \usepackage[inline]{enumitem} % IRkernel/repr support (it uses the enumerate* environment)
    \usepackage[normalem]{ulem} % ulem is needed to support strikethroughs (\sout)
                                % normalem makes italics be italics, not underlines
    \usepackage{soul}      % strikethrough (\st) support for pandoc >= 3.0.0
    \usepackage{mathrsfs}
    

    
    % Colors for the hyperref package
    \definecolor{urlcolor}{rgb}{0,.145,.698}
    \definecolor{linkcolor}{rgb}{.71,0.21,0.01}
    \definecolor{citecolor}{rgb}{.12,.54,.11}

    % ANSI colors
    \definecolor{ansi-black}{HTML}{3E424D}
    \definecolor{ansi-black-intense}{HTML}{282C36}
    \definecolor{ansi-red}{HTML}{E75C58}
    \definecolor{ansi-red-intense}{HTML}{B22B31}
    \definecolor{ansi-green}{HTML}{00A250}
    \definecolor{ansi-green-intense}{HTML}{007427}
    \definecolor{ansi-yellow}{HTML}{DDB62B}
    \definecolor{ansi-yellow-intense}{HTML}{B27D12}
    \definecolor{ansi-blue}{HTML}{208FFB}
    \definecolor{ansi-blue-intense}{HTML}{0065CA}
    \definecolor{ansi-magenta}{HTML}{D160C4}
    \definecolor{ansi-magenta-intense}{HTML}{A03196}
    \definecolor{ansi-cyan}{HTML}{60C6C8}
    \definecolor{ansi-cyan-intense}{HTML}{258F8F}
    \definecolor{ansi-white}{HTML}{C5C1B4}
    \definecolor{ansi-white-intense}{HTML}{A1A6B2}
    \definecolor{ansi-default-inverse-fg}{HTML}{FFFFFF}
    \definecolor{ansi-default-inverse-bg}{HTML}{000000}

    % common color for the border for error outputs.
    \definecolor{outerrorbackground}{HTML}{FFDFDF}

    % commands and environments needed by pandoc snippets
    % extracted from the output of `pandoc -s`
    \providecommand{\tightlist}{%
      \setlength{\itemsep}{0pt}\setlength{\parskip}{0pt}}
    \DefineVerbatimEnvironment{Highlighting}{Verbatim}{commandchars=\\\{\}}
    % Add ',fontsize=\small' for more characters per line
    \newenvironment{Shaded}{}{}
    \newcommand{\KeywordTok}[1]{\textcolor[rgb]{0.00,0.44,0.13}{\textbf{{#1}}}}
    \newcommand{\DataTypeTok}[1]{\textcolor[rgb]{0.56,0.13,0.00}{{#1}}}
    \newcommand{\DecValTok}[1]{\textcolor[rgb]{0.25,0.63,0.44}{{#1}}}
    \newcommand{\BaseNTok}[1]{\textcolor[rgb]{0.25,0.63,0.44}{{#1}}}
    \newcommand{\FloatTok}[1]{\textcolor[rgb]{0.25,0.63,0.44}{{#1}}}
    \newcommand{\CharTok}[1]{\textcolor[rgb]{0.25,0.44,0.63}{{#1}}}
    \newcommand{\StringTok}[1]{\textcolor[rgb]{0.25,0.44,0.63}{{#1}}}
    \newcommand{\CommentTok}[1]{\textcolor[rgb]{0.38,0.63,0.69}{\textit{{#1}}}}
    \newcommand{\OtherTok}[1]{\textcolor[rgb]{0.00,0.44,0.13}{{#1}}}
    \newcommand{\AlertTok}[1]{\textcolor[rgb]{1.00,0.00,0.00}{\textbf{{#1}}}}
    \newcommand{\FunctionTok}[1]{\textcolor[rgb]{0.02,0.16,0.49}{{#1}}}
    \newcommand{\RegionMarkerTok}[1]{{#1}}
    \newcommand{\ErrorTok}[1]{\textcolor[rgb]{1.00,0.00,0.00}{\textbf{{#1}}}}
    \newcommand{\NormalTok}[1]{{#1}}

    % Additional commands for more recent versions of Pandoc
    \newcommand{\ConstantTok}[1]{\textcolor[rgb]{0.53,0.00,0.00}{{#1}}}
    \newcommand{\SpecialCharTok}[1]{\textcolor[rgb]{0.25,0.44,0.63}{{#1}}}
    \newcommand{\VerbatimStringTok}[1]{\textcolor[rgb]{0.25,0.44,0.63}{{#1}}}
    \newcommand{\SpecialStringTok}[1]{\textcolor[rgb]{0.73,0.40,0.53}{{#1}}}
    \newcommand{\ImportTok}[1]{{#1}}
    \newcommand{\DocumentationTok}[1]{\textcolor[rgb]{0.73,0.13,0.13}{\textit{{#1}}}}
    \newcommand{\AnnotationTok}[1]{\textcolor[rgb]{0.38,0.63,0.69}{\textbf{\textit{{#1}}}}}
    \newcommand{\CommentVarTok}[1]{\textcolor[rgb]{0.38,0.63,0.69}{\textbf{\textit{{#1}}}}}
    \newcommand{\VariableTok}[1]{\textcolor[rgb]{0.10,0.09,0.49}{{#1}}}
    \newcommand{\ControlFlowTok}[1]{\textcolor[rgb]{0.00,0.44,0.13}{\textbf{{#1}}}}
    \newcommand{\OperatorTok}[1]{\textcolor[rgb]{0.40,0.40,0.40}{{#1}}}
    \newcommand{\BuiltInTok}[1]{{#1}}
    \newcommand{\ExtensionTok}[1]{{#1}}
    \newcommand{\PreprocessorTok}[1]{\textcolor[rgb]{0.74,0.48,0.00}{{#1}}}
    \newcommand{\AttributeTok}[1]{\textcolor[rgb]{0.49,0.56,0.16}{{#1}}}
    \newcommand{\InformationTok}[1]{\textcolor[rgb]{0.38,0.63,0.69}{\textbf{\textit{{#1}}}}}
    \newcommand{\WarningTok}[1]{\textcolor[rgb]{0.38,0.63,0.69}{\textbf{\textit{{#1}}}}}


    % Define a nice break command that doesn't care if a line doesn't already
    % exist.
    \def\br{\hspace*{\fill} \\* }
    % Math Jax compatibility definitions
    \def\gt{>}
    \def\lt{<}
    \let\Oldtex\TeX
    \let\Oldlatex\LaTeX
    \renewcommand{\TeX}{\textrm{\Oldtex}}
    \renewcommand{\LaTeX}{\textrm{\Oldlatex}}
    % Document parameters
    % Document title
    \title{agentic-memory-atomicity}
    
    
    
    
    
    
    
% Pygments definitions
\makeatletter
\def\PY@reset{\let\PY@it=\relax \let\PY@bf=\relax%
    \let\PY@ul=\relax \let\PY@tc=\relax%
    \let\PY@bc=\relax \let\PY@ff=\relax}
\def\PY@tok#1{\csname PY@tok@#1\endcsname}
\def\PY@toks#1+{\ifx\relax#1\empty\else%
    \PY@tok{#1}\expandafter\PY@toks\fi}
\def\PY@do#1{\PY@bc{\PY@tc{\PY@ul{%
    \PY@it{\PY@bf{\PY@ff{#1}}}}}}}
\def\PY#1#2{\PY@reset\PY@toks#1+\relax+\PY@do{#2}}

\@namedef{PY@tok@w}{\def\PY@tc##1{\textcolor[rgb]{0.73,0.73,0.73}{##1}}}
\@namedef{PY@tok@c}{\let\PY@it=\textit\def\PY@tc##1{\textcolor[rgb]{0.24,0.48,0.48}{##1}}}
\@namedef{PY@tok@cp}{\def\PY@tc##1{\textcolor[rgb]{0.61,0.40,0.00}{##1}}}
\@namedef{PY@tok@k}{\let\PY@bf=\textbf\def\PY@tc##1{\textcolor[rgb]{0.00,0.50,0.00}{##1}}}
\@namedef{PY@tok@kp}{\def\PY@tc##1{\textcolor[rgb]{0.00,0.50,0.00}{##1}}}
\@namedef{PY@tok@kt}{\def\PY@tc##1{\textcolor[rgb]{0.69,0.00,0.25}{##1}}}
\@namedef{PY@tok@o}{\def\PY@tc##1{\textcolor[rgb]{0.40,0.40,0.40}{##1}}}
\@namedef{PY@tok@ow}{\let\PY@bf=\textbf\def\PY@tc##1{\textcolor[rgb]{0.67,0.13,1.00}{##1}}}
\@namedef{PY@tok@nb}{\def\PY@tc##1{\textcolor[rgb]{0.00,0.50,0.00}{##1}}}
\@namedef{PY@tok@nf}{\def\PY@tc##1{\textcolor[rgb]{0.00,0.00,1.00}{##1}}}
\@namedef{PY@tok@nc}{\let\PY@bf=\textbf\def\PY@tc##1{\textcolor[rgb]{0.00,0.00,1.00}{##1}}}
\@namedef{PY@tok@nn}{\let\PY@bf=\textbf\def\PY@tc##1{\textcolor[rgb]{0.00,0.00,1.00}{##1}}}
\@namedef{PY@tok@ne}{\let\PY@bf=\textbf\def\PY@tc##1{\textcolor[rgb]{0.80,0.25,0.22}{##1}}}
\@namedef{PY@tok@nv}{\def\PY@tc##1{\textcolor[rgb]{0.10,0.09,0.49}{##1}}}
\@namedef{PY@tok@no}{\def\PY@tc##1{\textcolor[rgb]{0.53,0.00,0.00}{##1}}}
\@namedef{PY@tok@nl}{\def\PY@tc##1{\textcolor[rgb]{0.46,0.46,0.00}{##1}}}
\@namedef{PY@tok@ni}{\let\PY@bf=\textbf\def\PY@tc##1{\textcolor[rgb]{0.44,0.44,0.44}{##1}}}
\@namedef{PY@tok@na}{\def\PY@tc##1{\textcolor[rgb]{0.41,0.47,0.13}{##1}}}
\@namedef{PY@tok@nt}{\let\PY@bf=\textbf\def\PY@tc##1{\textcolor[rgb]{0.00,0.50,0.00}{##1}}}
\@namedef{PY@tok@nd}{\def\PY@tc##1{\textcolor[rgb]{0.67,0.13,1.00}{##1}}}
\@namedef{PY@tok@s}{\def\PY@tc##1{\textcolor[rgb]{0.73,0.13,0.13}{##1}}}
\@namedef{PY@tok@sd}{\let\PY@it=\textit\def\PY@tc##1{\textcolor[rgb]{0.73,0.13,0.13}{##1}}}
\@namedef{PY@tok@si}{\let\PY@bf=\textbf\def\PY@tc##1{\textcolor[rgb]{0.64,0.35,0.47}{##1}}}
\@namedef{PY@tok@se}{\let\PY@bf=\textbf\def\PY@tc##1{\textcolor[rgb]{0.67,0.36,0.12}{##1}}}
\@namedef{PY@tok@sr}{\def\PY@tc##1{\textcolor[rgb]{0.64,0.35,0.47}{##1}}}
\@namedef{PY@tok@ss}{\def\PY@tc##1{\textcolor[rgb]{0.10,0.09,0.49}{##1}}}
\@namedef{PY@tok@sx}{\def\PY@tc##1{\textcolor[rgb]{0.00,0.50,0.00}{##1}}}
\@namedef{PY@tok@m}{\def\PY@tc##1{\textcolor[rgb]{0.40,0.40,0.40}{##1}}}
\@namedef{PY@tok@gh}{\let\PY@bf=\textbf\def\PY@tc##1{\textcolor[rgb]{0.00,0.00,0.50}{##1}}}
\@namedef{PY@tok@gu}{\let\PY@bf=\textbf\def\PY@tc##1{\textcolor[rgb]{0.50,0.00,0.50}{##1}}}
\@namedef{PY@tok@gd}{\def\PY@tc##1{\textcolor[rgb]{0.63,0.00,0.00}{##1}}}
\@namedef{PY@tok@gi}{\def\PY@tc##1{\textcolor[rgb]{0.00,0.52,0.00}{##1}}}
\@namedef{PY@tok@gr}{\def\PY@tc##1{\textcolor[rgb]{0.89,0.00,0.00}{##1}}}
\@namedef{PY@tok@ge}{\let\PY@it=\textit}
\@namedef{PY@tok@gs}{\let\PY@bf=\textbf}
\@namedef{PY@tok@gp}{\let\PY@bf=\textbf\def\PY@tc##1{\textcolor[rgb]{0.00,0.00,0.50}{##1}}}
\@namedef{PY@tok@go}{\def\PY@tc##1{\textcolor[rgb]{0.44,0.44,0.44}{##1}}}
\@namedef{PY@tok@gt}{\def\PY@tc##1{\textcolor[rgb]{0.00,0.27,0.87}{##1}}}
\@namedef{PY@tok@err}{\def\PY@bc##1{{\setlength{\fboxsep}{\string -\fboxrule}\fcolorbox[rgb]{1.00,0.00,0.00}{1,1,1}{\strut ##1}}}}
\@namedef{PY@tok@kc}{\let\PY@bf=\textbf\def\PY@tc##1{\textcolor[rgb]{0.00,0.50,0.00}{##1}}}
\@namedef{PY@tok@kd}{\let\PY@bf=\textbf\def\PY@tc##1{\textcolor[rgb]{0.00,0.50,0.00}{##1}}}
\@namedef{PY@tok@kn}{\let\PY@bf=\textbf\def\PY@tc##1{\textcolor[rgb]{0.00,0.50,0.00}{##1}}}
\@namedef{PY@tok@kr}{\let\PY@bf=\textbf\def\PY@tc##1{\textcolor[rgb]{0.00,0.50,0.00}{##1}}}
\@namedef{PY@tok@bp}{\def\PY@tc##1{\textcolor[rgb]{0.00,0.50,0.00}{##1}}}
\@namedef{PY@tok@fm}{\def\PY@tc##1{\textcolor[rgb]{0.00,0.00,1.00}{##1}}}
\@namedef{PY@tok@vc}{\def\PY@tc##1{\textcolor[rgb]{0.10,0.09,0.49}{##1}}}
\@namedef{PY@tok@vg}{\def\PY@tc##1{\textcolor[rgb]{0.10,0.09,0.49}{##1}}}
\@namedef{PY@tok@vi}{\def\PY@tc##1{\textcolor[rgb]{0.10,0.09,0.49}{##1}}}
\@namedef{PY@tok@vm}{\def\PY@tc##1{\textcolor[rgb]{0.10,0.09,0.49}{##1}}}
\@namedef{PY@tok@sa}{\def\PY@tc##1{\textcolor[rgb]{0.73,0.13,0.13}{##1}}}
\@namedef{PY@tok@sb}{\def\PY@tc##1{\textcolor[rgb]{0.73,0.13,0.13}{##1}}}
\@namedef{PY@tok@sc}{\def\PY@tc##1{\textcolor[rgb]{0.73,0.13,0.13}{##1}}}
\@namedef{PY@tok@dl}{\def\PY@tc##1{\textcolor[rgb]{0.73,0.13,0.13}{##1}}}
\@namedef{PY@tok@s2}{\def\PY@tc##1{\textcolor[rgb]{0.73,0.13,0.13}{##1}}}
\@namedef{PY@tok@sh}{\def\PY@tc##1{\textcolor[rgb]{0.73,0.13,0.13}{##1}}}
\@namedef{PY@tok@s1}{\def\PY@tc##1{\textcolor[rgb]{0.73,0.13,0.13}{##1}}}
\@namedef{PY@tok@mb}{\def\PY@tc##1{\textcolor[rgb]{0.40,0.40,0.40}{##1}}}
\@namedef{PY@tok@mf}{\def\PY@tc##1{\textcolor[rgb]{0.40,0.40,0.40}{##1}}}
\@namedef{PY@tok@mh}{\def\PY@tc##1{\textcolor[rgb]{0.40,0.40,0.40}{##1}}}
\@namedef{PY@tok@mi}{\def\PY@tc##1{\textcolor[rgb]{0.40,0.40,0.40}{##1}}}
\@namedef{PY@tok@il}{\def\PY@tc##1{\textcolor[rgb]{0.40,0.40,0.40}{##1}}}
\@namedef{PY@tok@mo}{\def\PY@tc##1{\textcolor[rgb]{0.40,0.40,0.40}{##1}}}
\@namedef{PY@tok@ch}{\let\PY@it=\textit\def\PY@tc##1{\textcolor[rgb]{0.24,0.48,0.48}{##1}}}
\@namedef{PY@tok@cm}{\let\PY@it=\textit\def\PY@tc##1{\textcolor[rgb]{0.24,0.48,0.48}{##1}}}
\@namedef{PY@tok@cpf}{\let\PY@it=\textit\def\PY@tc##1{\textcolor[rgb]{0.24,0.48,0.48}{##1}}}
\@namedef{PY@tok@c1}{\let\PY@it=\textit\def\PY@tc##1{\textcolor[rgb]{0.24,0.48,0.48}{##1}}}
\@namedef{PY@tok@cs}{\let\PY@it=\textit\def\PY@tc##1{\textcolor[rgb]{0.24,0.48,0.48}{##1}}}

\def\PYZbs{\char`\\}
\def\PYZus{\char`\_}
\def\PYZob{\char`\{}
\def\PYZcb{\char`\}}
\def\PYZca{\char`\^}
\def\PYZam{\char`\&}
\def\PYZlt{\char`\<}
\def\PYZgt{\char`\>}
\def\PYZsh{\char`\#}
\def\PYZpc{\char`\%}
\def\PYZdl{\char`\$}
\def\PYZhy{\char`\-}
\def\PYZsq{\char`\'}
\def\PYZdq{\char`\"}
\def\PYZti{\char`\~}
% for compatibility with earlier versions
\def\PYZat{@}
\def\PYZlb{[}
\def\PYZrb{]}
\makeatother


    % For linebreaks inside Verbatim environment from package fancyvrb.
    \makeatletter
        \newbox\Wrappedcontinuationbox
        \newbox\Wrappedvisiblespacebox
        \newcommand*\Wrappedvisiblespace {\textcolor{red}{\textvisiblespace}}
        \newcommand*\Wrappedcontinuationsymbol {\textcolor{red}{\llap{\tiny$\m@th\hookrightarrow$}}}
        \newcommand*\Wrappedcontinuationindent {3ex }
        \newcommand*\Wrappedafterbreak {\kern\Wrappedcontinuationindent\copy\Wrappedcontinuationbox}
        % Take advantage of the already applied Pygments mark-up to insert
        % potential linebreaks for TeX processing.
        %        {, <, #, %, $, ' and ": go to next line.
        %        _, }, ^, &, >, - and ~: stay at end of broken line.
        % Use of \textquotesingle for straight quote.
        \newcommand*\Wrappedbreaksatspecials {%
            \def\PYGZus{\discretionary{\char`\_}{\Wrappedafterbreak}{\char`\_}}%
            \def\PYGZob{\discretionary{}{\Wrappedafterbreak\char`\{}{\char`\{}}%
            \def\PYGZcb{\discretionary{\char`\}}{\Wrappedafterbreak}{\char`\}}}%
            \def\PYGZca{\discretionary{\char`\^}{\Wrappedafterbreak}{\char`\^}}%
            \def\PYGZam{\discretionary{\char`\&}{\Wrappedafterbreak}{\char`\&}}%
            \def\PYGZlt{\discretionary{}{\Wrappedafterbreak\char`\<}{\char`\<}}%
            \def\PYGZgt{\discretionary{\char`\>}{\Wrappedafterbreak}{\char`\>}}%
            \def\PYGZsh{\discretionary{}{\Wrappedafterbreak\char`\#}{\char`\#}}%
            \def\PYGZpc{\discretionary{}{\Wrappedafterbreak\char`\%}{\char`\%}}%
            \def\PYGZdl{\discretionary{}{\Wrappedafterbreak\char`\$}{\char`\$}}%
            \def\PYGZhy{\discretionary{\char`\-}{\Wrappedafterbreak}{\char`\-}}%
            \def\PYGZsq{\discretionary{}{\Wrappedafterbreak\textquotesingle}{\textquotesingle}}%
            \def\PYGZdq{\discretionary{}{\Wrappedafterbreak\char`\"}{\char`\"}}%
            \def\PYGZti{\discretionary{\char`\~}{\Wrappedafterbreak}{\char`\~}}%
        }
        % Some characters . , ; ? ! / are not pygmentized.
        % This macro makes them "active" and they will insert potential linebreaks
        \newcommand*\Wrappedbreaksatpunct {%
            \lccode`\~`\.\lowercase{\def~}{\discretionary{\hbox{\char`\.}}{\Wrappedafterbreak}{\hbox{\char`\.}}}%
            \lccode`\~`\,\lowercase{\def~}{\discretionary{\hbox{\char`\,}}{\Wrappedafterbreak}{\hbox{\char`\,}}}%
            \lccode`\~`\;\lowercase{\def~}{\discretionary{\hbox{\char`\;}}{\Wrappedafterbreak}{\hbox{\char`\;}}}%
            \lccode`\~`\:\lowercase{\def~}{\discretionary{\hbox{\char`\:}}{\Wrappedafterbreak}{\hbox{\char`\:}}}%
            \lccode`\~`\?\lowercase{\def~}{\discretionary{\hbox{\char`\?}}{\Wrappedafterbreak}{\hbox{\char`\?}}}%
            \lccode`\~`\!\lowercase{\def~}{\discretionary{\hbox{\char`\!}}{\Wrappedafterbreak}{\hbox{\char`\!}}}%
            \lccode`\~`\/\lowercase{\def~}{\discretionary{\hbox{\char`\/}}{\Wrappedafterbreak}{\hbox{\char`\/}}}%
            \catcode`\.\active
            \catcode`\,\active
            \catcode`\;\active
            \catcode`\:\active
            \catcode`\?\active
            \catcode`\!\active
            \catcode`\/\active
            \lccode`\~`\~
        }
    \makeatother

    \let\OriginalVerbatim=\Verbatim
    \makeatletter
    \renewcommand{\Verbatim}[1][1]{%
        %\parskip\z@skip
        \sbox\Wrappedcontinuationbox {\Wrappedcontinuationsymbol}%
        \sbox\Wrappedvisiblespacebox {\FV@SetupFont\Wrappedvisiblespace}%
        \def\FancyVerbFormatLine ##1{\hsize\linewidth
            \vtop{\raggedright\hyphenpenalty\z@\exhyphenpenalty\z@
                \doublehyphendemerits\z@\finalhyphendemerits\z@
                \strut ##1\strut}%
        }%
        % If the linebreak is at a space, the latter will be displayed as visible
        % space at end of first line, and a continuation symbol starts next line.
        % Stretch/shrink are however usually zero for typewriter font.
        \def\FV@Space {%
            \nobreak\hskip\z@ plus\fontdimen3\font minus\fontdimen4\font
            \discretionary{\copy\Wrappedvisiblespacebox}{\Wrappedafterbreak}
            {\kern\fontdimen2\font}%
        }%

        % Allow breaks at special characters using \PYG... macros.
        \Wrappedbreaksatspecials
        % Breaks at punctuation characters . , ; ? ! and / need catcode=\active
        \OriginalVerbatim[#1,codes*=\Wrappedbreaksatpunct]%
    }
    \makeatother

    % Exact colors from NB
    \definecolor{incolor}{HTML}{303F9F}
    \definecolor{outcolor}{HTML}{D84315}
    \definecolor{cellborder}{HTML}{CFCFCF}
    \definecolor{cellbackground}{HTML}{F7F7F7}

    % prompt
    \makeatletter
    \newcommand{\boxspacing}{\kern\kvtcb@left@rule\kern\kvtcb@boxsep}
    \makeatother
    \newcommand{\prompt}[4]{
        {\ttfamily\llap{{\color{#2}[#3]:\hspace{3pt}#4}}\vspace{-\baselineskip}}
    }
    

    
    % Prevent overflowing lines due to hard-to-break entities
    \sloppy
    % Setup hyperref package
    \hypersetup{
      breaklinks=true,  % so long urls are correctly broken across lines
      colorlinks=true,
      urlcolor=urlcolor,
      linkcolor=linkcolor,
      citecolor=citecolor,
      }
    % Slightly bigger margins than the latex defaults
    
    \geometry{verbose,tmargin=1in,bmargin=1in,lmargin=1in,rmargin=1in}
    
    

\begin{document}
    
    \maketitle
    
    

    
    \section{Agentic Memory Atomicity:}\label{agentic-memory-atomicity}

\section{A Transactional Memory Architecture for Ephemeral AI Agent
Systems}\label{a-transactional-memory-architecture-for-ephemeral-ai-agent-systems}

\textbf{Loïc Mancino} · \textbf{Jérôme Chincarini}

\emph{Synergix Labs --- 1 year of R\&D on
\href{https://wrai.th}{wrai.th}}

\begin{center}\rule{0.5\linewidth}{0.5pt}\end{center}

\textbf{Abstract} --- We present \emph{Agentic Memory Atomicity}, an
architecture where AI agents are stateless ephemeral processes that read
from and commit to a persistent, transactional memory store. Each agent
spawns, loads inherited knowledge, performs work, commits new knowledge
atomically, and terminates. The next spawn inherits everything. We
formalize the state model, demonstrate monotonic knowledge growth under
compaction, and show sub-linear cost scaling for multi-agent teams
sharing project-scoped memory. Our implementation (relay OS) uses SQLite
with FTS5, MCP protocol, and Go, with mechanisms for conflict
resolution, budget pruning, temporal garbage collection, and
priority-based context injection. Simulations show convergence of
marginal learning, decreasing cost per unit of work, and robustness to
LLM noise.

\emph{Keywords: multi-agent systems, persistent memory, transactional
state, ephemeral compute, LLM agents}

    \subsection{1. Introduction}\label{introduction}

Current LLM-based agent systems face a fundamental tension: agents need
persistent knowledge to be effective, but long-running processes suffer
from context window exhaustion, idle cost, and crash fragility.

We propose \textbf{Agentic Memory Atomicity} --- a design where:

\begin{enumerate}
\def\labelenumi{\arabic{enumi}.}
\tightlist
\item
  Agents are \textbf{ephemeral}: they spawn, work, and die
\item
  Knowledge lives in a \textbf{persistent, transactional store} outside
  the agent
\item
  Each spawn \textbf{inherits all prior knowledge} via a projection
  function
\item
  Memory commits are \textbf{atomic}: success writes \(\Delta M\), crash
  preserves \(M_t\)
\end{enumerate}

This paper formalizes the model, identifies threats (LLM noise, memory
explosion, multi-agent conflicts, temporal staleness), and demonstrates
how implemented mechanisms address each threat.

    \subsection{2. Background \& Related
Work}\label{background-related-work}

\subsubsection{Existing Approaches}\label{existing-approaches}

\begin{longtable}[]{@{}
  >{\raggedright\arraybackslash}p{(\linewidth - 6\tabcolsep) * \real{0.1633}}
  >{\raggedright\arraybackslash}p{(\linewidth - 6\tabcolsep) * \real{0.2653}}
  >{\raggedright\arraybackslash}p{(\linewidth - 6\tabcolsep) * \real{0.3265}}
  >{\raggedright\arraybackslash}p{(\linewidth - 6\tabcolsep) * \real{0.2449}}@{}}
\toprule\noalign{}
\begin{minipage}[b]{\linewidth}\raggedright
System
\end{minipage} & \begin{minipage}[b]{\linewidth}\raggedright
Memory Model
\end{minipage} & \begin{minipage}[b]{\linewidth}\raggedright
Agent Lifecycle
\end{minipage} & \begin{minipage}[b]{\linewidth}\raggedright
Limitation
\end{minipage} \\
\midrule\noalign{}
\endhead
\bottomrule\noalign{}
\endlastfoot
\textbf{LangGraph} & Graph state (in-memory) & Persistent process &
Context exhaustion, no crash recovery \\
\textbf{CrewAI} & Shared task queue & Persistent threads & No
transactional memory, no inheritance \\
\textbf{AutoGen} & Conversation history & Session-bound & Lost on
termination \\
\textbf{Relay OS} (ours) & Transactional DB (SQLite) & Ephemeral spawns
& --- \\
\end{longtable}

\subsubsection{Key Differentiator}\label{key-differentiator}

Existing frameworks treat memory as a side-effect of agent execution. We
treat it as the \textbf{primary artifact} --- agents are disposable
compute, memory is the durable intelligence.

\[\text{Intelligence} = f(M_t), \quad \text{Agent} = \text{ephemeral compute over } M_t\]

    \subsection{3. Agent State Model}\label{agent-state-model}

\subsubsection{Definition 1 --- Agent
State}\label{definition-1-agent-state}

\[S_t(A) = \{ I, M_t, V, T_t, \Omega_t, R \}\]

\begin{longtable}[]{@{}lll@{}}
\toprule\noalign{}
Symbol & Name & Nature \\
\midrule\noalign{}
\endhead
\bottomrule\noalign{}
\endlastfoot
\(I\) & Identity (profile, role, reports\_to) & Static \\
\(M_t\) & Memories at time \(t\) --- 3 scopes & Monotonically growing \\
\(V\) & Vault (conventions, docs) & Quasi-static \\
\(T_t\) & Assigned tasks & Variable \\
\(\Omega_t\) & Inbox (unread messages) & Ephemeral, TTL-bound \\
\(R\) & Rules (boot, cycle, learning) & Static \\
\end{longtable}

\subsubsection{Definition 2 --- Memory
Scopes}\label{definition-2-memory-scopes}

\[M_t = M_t^{\text{agent}} \cup M_t^{\text{project}} \cup M_t^{\text{global}}\]

\begin{longtable}[]{@{}lll@{}}
\toprule\noalign{}
Scope & Visibility & Cascade Priority \\
\midrule\noalign{}
\endhead
\bottomrule\noalign{}
\endlastfoot
\texttt{agent} & Private to owner & 1 (checked first) \\
\texttt{project} & All agents in project & 2 \\
\texttt{global} & All agents everywhere & 3 (fallback) \\
\end{longtable}

Lookup: \texttt{agent\ →\ project\ →\ global} (first match wins).

\subsubsection{Definition 3 --- Memory
Layers}\label{definition-3-memory-layers}

\[\text{constraints} > \text{behavior} > \text{context}\]

\begin{itemize}
\tightlist
\item
  \texttt{constraints}: hard rules (never overridden)
\item
  \texttt{behavior}: adaptive defaults
\item
  \texttt{context}: ephemeral, session-specific
\end{itemize}

    \begin{center}
    \adjustimage{max size={0.9\linewidth}{0.9\paperheight}}{agentic-memory-atomicity_files/agentic-memory-atomicity_5_0.png}
    \end{center}
    { \hspace*{\fill} \\}
    
    \subsection{4. Ephemeral Agent
Principle}\label{ephemeral-agent-principle}

\subsubsection{Principle}\label{principle}

Agents are \textbf{stateless compute processes}. They carry no internal
state between invocations.

\[\text{Agent} = f(M_t) \rightarrow (W, \Delta M)\]

The intelligence of the system resides entirely in \(M_t\), the
persistent memory store. An agent is a pure function that:

\begin{enumerate}
\def\labelenumi{\arabic{enumi}.}
\tightlist
\item
  \textbf{Spawns} --- process created by scheduler (cron) or dispatch
\item
  \textbf{Loads memory} --- \texttt{BuildSpawnContext()} projects
  \(M_t\) into the agent's context
\item
  \textbf{Computes} --- executes work \(W\) (code, messages, task
  updates)
\item
  \textbf{Commits} \(\Delta M\) --- new knowledge persisted atomically
\item
  \textbf{Terminates} --- process dies, zero residual cost
\end{enumerate}

\subsubsection{Contrast with Persistent
Agents}\label{contrast-with-persistent-agents}

\begin{longtable}[]{@{}
  >{\raggedright\arraybackslash}p{(\linewidth - 4\tabcolsep) * \real{0.2326}}
  >{\raggedright\arraybackslash}p{(\linewidth - 4\tabcolsep) * \real{0.3953}}
  >{\raggedright\arraybackslash}p{(\linewidth - 4\tabcolsep) * \real{0.3721}}@{}}
\toprule\noalign{}
\begin{minipage}[b]{\linewidth}\raggedright
Property
\end{minipage} & \begin{minipage}[b]{\linewidth}\raggedright
Persistent Agent
\end{minipage} & \begin{minipage}[b]{\linewidth}\raggedright
Ephemeral Agent
\end{minipage} \\
\midrule\noalign{}
\endhead
\bottomrule\noalign{}
\endlastfoot
Context window & Exhausts → loses early context & Always fresh \\
Crash recovery & Corrupted state possible & \(M_t\) intact \\
Idle cost & Tokens burned (keepalive) & Zero \\
Scalability & 1 process = 1 agent & Pool of \(N\) spawns \\
Knowledge & Limited to session & Cumulative \(\infty\) \\
Formalism & \(S_t = f(S_{t-1}, \text{input})\) (stateful) &
\(c : S_t \to (S_{t+1}, W)\) (functional) \\
\end{longtable}

    \begin{center}
    \adjustimage{max size={0.9\linewidth}{0.9\paperheight}}{agentic-memory-atomicity_files/agentic-memory-atomicity_7_0.png}
    \end{center}
    { \hspace*{\fill} \\}
    
    \subsection{5. Atomic Memory Commit}\label{atomic-memory-commit}

\subsubsection{Definition 4 --- Cycle
Transition}\label{definition-4-cycle-transition}

\[c : S_t \rightarrow (S_{t+1}, W)\]

Where \(W\) = work produced. The state transition decomposes:

\begin{itemize}
\tightlist
\item
  \(M_{t+1} = \text{compact}(M_t \cup \Delta M_{\text{valid}})\) ---
  memory updated (filtered + compacted)
\item
  \(\Omega_{t+1}' = \Omega_t \cup \Delta \Omega\) --- messages emitted
  to other agents
\item
  \(T_{t+1}\) --- tasks updated (claimed, completed, blocked)
\end{itemize}

\subsubsection{Definition 5 --- Atomicity
Property}\label{definition-5-atomicity-property}

\[M_{t+1} = \begin{cases} \text{compact}(M_t \cup \Delta M_{\text{valid}}) & \text{if cycle succeeds (exit 0)} \\ M_t & \text{if cycle crashes} \end{cases}\]

No intermediate corrupted state. Transaction commit or rollback.

\subsubsection{Definition 6 --- Memory Validation
Filter}\label{definition-6-memory-validation-filter}

\[\Delta M_{\text{valid}} = \text{filter}(\Delta M, M_t)\]

The filter enforces: - \textbf{Conflict detection}: two agents writing
same key → both preserved until explicit resolution - \textbf{Layer
hierarchy}:
\texttt{constraints\ \textgreater{}\ behavior\ \textgreater{}\ context}
--- a \texttt{behavior} write cannot override \texttt{constraints} -
\textbf{Scope isolation}: \texttt{agent} (private) / \texttt{project}
(shared) / \texttt{global} (cross-project) - \textbf{Deduplication}:
same key+value → timestamp update only, no new version

\subsubsection{Definition 7 --- Memory
Compaction}\label{definition-7-memory-compaction}

\[|M_{\text{boot}}| \leq B_{\max}\]

The raw memory grows monotonically, but the \textbf{projected view} at
boot is bounded by a byte budget.

    \begin{center}
    \adjustimage{max size={0.9\linewidth}{0.9\paperheight}}{agentic-memory-atomicity_files/agentic-memory-atomicity_9_0.png}
    \end{center}
    { \hspace*{\fill} \\}
    
    \subsection{6. Knowledge Evolution}\label{knowledge-evolution}

\subsubsection{Spawn Chain Inheritance}\label{spawn-chain-inheritance}

For two successive spawns \(s_1\) and \(s_2\) of the same profile:

\[M(s_2) \supseteq M(s_1)\]

Agent \(s_2\) \textbf{knows everything} \(s_1\) learned. This is not
shared memory at runtime --- it is \textbf{boot-time inheritance} via
\texttt{BuildSpawnContext()}.

\subsubsection{Inter-Agent Memory
Elimination}\label{inter-agent-memory-elimination}

An agent can supersede another agent's knowledge:

\[\text{set\_memory}(\text{key}, v_{\text{new}}, \text{scope:project}, \text{upsert:true})\]

Archives the old version (field \texttt{supersedes}), increments
\texttt{version}. Complete traceability, never hard-delete.

\[v_1 \xrightarrow{\text{supersedes}} v_2 \xrightarrow{\text{supersedes}} v_3\]

\subsubsection{Self-Learning Correction}\label{self-learning-correction}

A Backend agent can correct a Frontend agent's project memory if it
discovers it is inaccurate. The \texttt{supersedes} chain preserves full
provenance.

    \begin{center}
    \adjustimage{max size={0.9\linewidth}{0.9\paperheight}}{agentic-memory-atomicity_files/agentic-memory-atomicity_11_0.png}
    \end{center}
    { \hspace*{\fill} \\}
    
    \begin{Verbatim}[commandchars=\\\{\}]
After 200 cycles:
  CTO: 338 items (30 private, 308 shared)
  Back-1: 338 items (30 private, 308 shared)
  Back-2: 336 items (28 private, 308 shared)
  Frontend: 338 items (30 private, 308 shared)
  DevOps: 338 items (30 private, 308 shared)
  QA: 338 items (30 private, 308 shared)
    \end{Verbatim}

    \begin{center}
    \adjustimage{max size={0.9\linewidth}{0.9\paperheight}}{agentic-memory-atomicity_files/agentic-memory-atomicity_12_0.png}
    \end{center}
    { \hspace*{\fill} \\}
    
    \begin{Verbatim}[commandchars=\\\{\}]
After 24h:
  Persistent: 19 items, 24.6M tokens, 96 context resets
  Ephemeral (1): 500 items, 6.6M tokens
  Ephemeral (6): 500 items, 17.7M tokens
  Multi-agent advantage: 26.3x knowledge, 28\% cost reduction vs persistent
    \end{Verbatim}

    \subsection{7. Convergence Conjecture}\label{convergence-conjecture}

\subsubsection{Conjecture: Knowledge Saturation
Hypothesis}\label{conjecture-knowledge-saturation-hypothesis}

\[\lim_{n \to \infty} |\Delta M_{\text{valid},n}| \to 0 \quad \text{(finite domain, active filter)}\]

\textbf{This is not a theorem.} It is an empirical hypothesis that holds
under three conditions:

\begin{enumerate}
\def\labelenumi{\arabic{enumi}.}
\tightlist
\item
  The knowledge domain is finite (bounded by the project)
\item
  The filter \(\text{filter}(\Delta M, M_t)\) eliminates duplicates and
  contradictions
\item
  The agent learns effectively (no hallucination loops)
\end{enumerate}

In practice, LLMs produce noise (repetitions, trivialities,
hallucinations). The observed convergence is due to the \textbf{filter +
compaction}, not an intrinsic property of learning.

\subsubsection{Corollary: Marginal Cost
Decrease}\label{corollary-marginal-cost-decrease}

\[C_t = \underbrace{C_{\text{boot}}}_{\text{fixed}} + \underbrace{C_{\text{context}}(|M_t^{\text{visible}}|)}_{\text{bounded by } B_{\max}} + \underbrace{C_{\text{work}}}_{\text{constant}} + \underbrace{C_{\text{learn}}(|\Delta M_t|)}_{\text{decreasing}}\]

As \(|\Delta M_t| \to 0\), the cost converges to
\(C_{\text{boot}} + C_{\text{context}} + C_{\text{work}}\) --- the
learning overhead vanishes.

    \begin{center}
    \adjustimage{max size={0.9\linewidth}{0.9\paperheight}}{agentic-memory-atomicity_files/agentic-memory-atomicity_14_0.png}
    \end{center}
    { \hspace*{\fill} \\}
    
    \begin{Verbatim}[commandchars=\\\{\}]
Boot tokens/day: 52.8M → 6.3M (88\% reduction)
Total tokens/day: 64.6M → 18.1M (72\% reduction)
Sonnet: \$5,815/mo → \$1,631/mo
Opus:   \$29,077/mo → \$8,156/mo
    \end{Verbatim}

    \subsection{8. Memory Validation \& Conflict
Resolution}\label{memory-validation-conflict-resolution}

\subsubsection{Validation Pipeline}\label{validation-pipeline}

Every memory write passes through a multi-stage filter before
persistence:

\[\Delta M \xrightarrow{\text{dedup}} \xrightarrow{\text{layer check}} \xrightarrow{\text{scope isolation}} \xrightarrow{\text{conflict detection}} \xrightarrow{\text{commit}} M_{t+1}\]

\subsubsection{Write Modes}\label{write-modes}

\begin{longtable}[]{@{}
  >{\raggedright\arraybackslash}p{(\linewidth - 4\tabcolsep) * \real{0.2400}}
  >{\raggedright\arraybackslash}p{(\linewidth - 4\tabcolsep) * \real{0.3600}}
  >{\raggedright\arraybackslash}p{(\linewidth - 4\tabcolsep) * \real{0.4000}}@{}}
\toprule\noalign{}
\begin{minipage}[b]{\linewidth}\raggedright
Mode
\end{minipage} & \begin{minipage}[b]{\linewidth}\raggedright
\texttt{upsert}
\end{minipage} & \begin{minipage}[b]{\linewidth}\raggedright
Behavior
\end{minipage} \\
\midrule\noalign{}
\endhead
\bottomrule\noalign{}
\endlastfoot
Silent overwrite & \texttt{true} (default) & Archive old →
\texttt{supersedes} chain → new version \\
Conflict detection & \texttt{false} & Preserve both →
\texttt{conflict\_with} cross-ref → await resolution \\
\end{longtable}

\subsubsection{Resolution}\label{resolution}

\texttt{resolve\_conflict(key,\ chosen\_value,\ scope)}: finds all
active versions, archives losers
(\texttt{archived\_by\ =\ "conflict\_resolution"}), clears conflict flag
on winner.

\subsubsection{Budget Pruning (Utility
Scoring)}\label{budget-pruning-utility-scoring}

\[U(m) = 0.7 \times \text{priority}(m) + 0.2 \times \text{Jaccard}(\text{tags}_m, \text{tags}_{\text{agent}}) + 0.1 \times \text{freshness}(m)\]

Where: -
\(\text{priority}: \text{P0}=1.0, \text{P1}\approx0.67, \text{P2}\approx0.33, \text{P3}=0\)
- \(\text{Jaccard}(A, B) = |A \cap B| / |A \cup B|\) -
\(\text{freshness}(m) = 1 / (1 + \text{age}(m) / 3600)\)

\textbf{P0 (interrupt) always bypasses the budget.} Remaining messages
scored and selected greedily.

    \begin{center}
    \adjustimage{max size={0.9\linewidth}{0.9\paperheight}}{agentic-memory-atomicity_files/agentic-memory-atomicity_16_0.png}
    \end{center}
    { \hspace*{\fill} \\}
    
    \begin{center}
    \adjustimage{max size={0.9\linewidth}{0.9\paperheight}}{agentic-memory-atomicity_files/agentic-memory-atomicity_17_0.png}
    \end{center}
    { \hspace*{\fill} \\}
    
    \begin{Verbatim}[commandchars=\\\{\}]
P0: 3/3 delivered (100\% — ALWAYS bypass)
P3 delivered: 0/12
    \end{Verbatim}

    \subsection{9. Implementation}\label{implementation}

\subsubsection{Stack}\label{stack}

\begin{longtable}[]{@{}
  >{\raggedright\arraybackslash}p{(\linewidth - 4\tabcolsep) * \real{0.3667}}
  >{\raggedright\arraybackslash}p{(\linewidth - 4\tabcolsep) * \real{0.3667}}
  >{\raggedright\arraybackslash}p{(\linewidth - 4\tabcolsep) * \real{0.2667}}@{}}
\toprule\noalign{}
\begin{minipage}[b]{\linewidth}\raggedright
Component
\end{minipage} & \begin{minipage}[b]{\linewidth}\raggedright
Technology
\end{minipage} & \begin{minipage}[b]{\linewidth}\raggedright
Purpose
\end{minipage} \\
\midrule\noalign{}
\endhead
\bottomrule\noalign{}
\endlastfoot
Relay server & Go & MCP server, HTTP API, SSE streaming \\
Storage & SQLite + FTS5 & Transactional memory, full-text search \\
Agent runtime & Claude Code (headless) & LLM compute via MCP \\
Scheduling & Cron expressions & Agent lifecycle management \\
Protocol & MCP (Model Context Protocol) & Tool discovery, push
notifications \\
\end{longtable}

\subsubsection{Key Functions}\label{key-functions}

\begin{longtable}[]{@{}
  >{\raggedright\arraybackslash}p{(\linewidth - 4\tabcolsep) * \real{0.4545}}
  >{\raggedright\arraybackslash}p{(\linewidth - 4\tabcolsep) * \real{0.2727}}
  >{\raggedright\arraybackslash}p{(\linewidth - 4\tabcolsep) * \real{0.2727}}@{}}
\toprule\noalign{}
\begin{minipage}[b]{\linewidth}\raggedright
Function
\end{minipage} & \begin{minipage}[b]{\linewidth}\raggedright
File
\end{minipage} & \begin{minipage}[b]{\linewidth}\raggedright
Role
\end{minipage} \\
\midrule\noalign{}
\endhead
\bottomrule\noalign{}
\endlastfoot
\texttt{BuildSpawnContext()} & \texttt{spawn/prompt.go} & Projection:
\(M_t \to S_t\) \\
\texttt{SetMemory()} & \texttt{db/memories.go} & Commit:
\(\Delta M \to M_{t+1}\) (upsert or conflict) \\
\texttt{ResolveConflict()} & \texttt{db/memories.go} & Conflict
resolution: pick winner, archive losers \\
\texttt{applyBudget()} & \texttt{relay/budget.go} & Budget pruning:
\(U(m) = 0.7P + 0.2J + 0.1F\) \\
\texttt{StartCleanup()} & \texttt{relay/cleanup.go} & Temporal GC: TTL,
staleness, ACK escalation \\
\texttt{FormatPrompt()} & \texttt{spawn/prompt.go} & \(S_t \to\) prompt
string (agent opens its eyes) \\
\end{longtable}

\subsubsection{Temporal Mechanisms}\label{temporal-mechanisms}

\begin{longtable}[]{@{}
  >{\raggedright\arraybackslash}p{(\linewidth - 4\tabcolsep) * \real{0.3793}}
  >{\raggedright\arraybackslash}p{(\linewidth - 4\tabcolsep) * \real{0.3448}}
  >{\raggedright\arraybackslash}p{(\linewidth - 4\tabcolsep) * \real{0.2759}}@{}}
\toprule\noalign{}
\begin{minipage}[b]{\linewidth}\raggedright
Mechanism
\end{minipage} & \begin{minipage}[b]{\linewidth}\raggedright
Interval
\end{minipage} & \begin{minipage}[b]{\linewidth}\raggedright
Timeout
\end{minipage} \\
\midrule\noalign{}
\endhead
\bottomrule\noalign{}
\endlastfoot
Message TTL & On send & Default 4h, \texttt{0} = immortal \\
Delivery state machine &
\texttt{queued\ →\ surfaced\ →\ ack\ →\ expired} & Cascades from message
TTL \\
ACK escalation & Every 5min & 15min → notify, 45min → escalate \\
Agent staleness & Every 5min & 30min inactivity → \texttt{inactive} \\
File lock TTL & Every 5min & Default 30min \\
Ghost reaping & Every 30s & Zombie spawns → \texttt{dead} \\
DB optimization & Every 5min & \texttt{PRAGMA\ optimize} + WAL
checkpoint \\
\end{longtable}

    \subsection{10. Experiments}\label{experiments}

We present three simulation experiments validating the core claims.

    \begin{center}
    \adjustimage{max size={0.9\linewidth}{0.9\paperheight}}{agentic-memory-atomicity_files/agentic-memory-atomicity_20_0.png}
    \end{center}
    { \hspace*{\fill} \\}
    
    \begin{center}
    \adjustimage{max size={0.9\linewidth}{0.9\paperheight}}{agentic-memory-atomicity_files/agentic-memory-atomicity_21_0.png}
    \end{center}
    { \hspace*{\fill} \\}
    
    \begin{Verbatim}[commandchars=\\\{\}]
Tasks: 5 dispatched, 5 acked, 0 escalated
    \end{Verbatim}

    \begin{center}
    \adjustimage{max size={0.9\linewidth}{0.9\paperheight}}{agentic-memory-atomicity_files/agentic-memory-atomicity_22_0.png}
    \end{center}
    { \hspace*{\fill} \\}
    
    \begin{Verbatim}[commandchars=\\\{\}]
Total: 14 conflicts, 1 resolutions, 23 supersedes
Active conflicts at end: 12
    \end{Verbatim}

    \subsection{11. Discussion \& Conclusion}\label{discussion-conclusion}

\subsubsection{Axioms}\label{axioms}

\begin{longtable}[]{@{}
  >{\raggedright\arraybackslash}p{(\linewidth - 4\tabcolsep) * \real{0.1500}}
  >{\raggedright\arraybackslash}p{(\linewidth - 4\tabcolsep) * \real{0.3500}}
  >{\raggedright\arraybackslash}p{(\linewidth - 4\tabcolsep) * \real{0.5000}}@{}}
\toprule\noalign{}
\begin{minipage}[b]{\linewidth}\raggedright
\#
\end{minipage} & \begin{minipage}[b]{\linewidth}\raggedright
Axiom
\end{minipage} & \begin{minipage}[b]{\linewidth}\raggedright
Statement
\end{minipage} \\
\midrule\noalign{}
\endhead
\bottomrule\noalign{}
\endlastfoot
A1 & Monotonicity & \(M_{t+1} \supseteq M_t\) modulo compaction \\
A2 & Atomicity & Cycle produces \((\Delta M, \Delta\Omega, W)\) or
nothing \\
A3 & Inheritance & \(M(s_{n+1}) \supseteq M(s_n)\) (subject to
\(B_{max}\)) \\
A4 & Saturation & \(\lim |\Delta M_n| \to 0\) (conjecture, finite
domain) \\
A5 & Zero idle & Terminated agent = 0 cost \\
A6 & Scopes & \(M = M^{agent} \cup M^{project} \cup M^{global}\),
cascade lookup \\
A7 & Layers & constraints \textgreater{} behavior \textgreater{} context
(never overridden) \\
\end{longtable}

\subsubsection{Implemented Mechanisms}\label{implemented-mechanisms}

\begin{longtable}[]{@{}
  >{\raggedright\arraybackslash}p{(\linewidth - 4\tabcolsep) * \real{0.1304}}
  >{\raggedright\arraybackslash}p{(\linewidth - 4\tabcolsep) * \real{0.4348}}
  >{\raggedright\arraybackslash}p{(\linewidth - 4\tabcolsep) * \real{0.4348}}@{}}
\toprule\noalign{}
\begin{minipage}[b]{\linewidth}\raggedright
\#
\end{minipage} & \begin{minipage}[b]{\linewidth}\raggedright
Mechanism
\end{minipage} & \begin{minipage}[b]{\linewidth}\raggedright
Addresses
\end{minipage} \\
\midrule\noalign{}
\endhead
\bottomrule\noalign{}
\endlastfoot
R1 & Budget pruning (\(U = 0.7P + 0.2J + 0.1F\), P0 bypass) & LLM noise,
boot cost explosion \\
R2 & Temporal GC (TTL, staleness, ACK escalation, ghost reaping) & Stale
messages, zombie agents, orphan tasks \\
R3 & Conflict resolution (upsert/conflict modes,
\texttt{resolve\_conflict}, autonomous) & Multi-agent divergence \\
R4 & Tool security
(\(\text{tools}_{eff} = \text{tools}_{cycle} \cap \text{tools}_{profile}\))
& Privilege escalation \\
R5 & Versioning (supersedes chain, soft-delete, archive) & Knowledge
provenance, audit trail \\
\end{longtable}

\subsubsection{Key Results}\label{key-results}

\begin{longtable}[]{@{}
  >{\raggedright\arraybackslash}p{(\linewidth - 2\tabcolsep) * \real{0.4118}}
  >{\raggedright\arraybackslash}p{(\linewidth - 2\tabcolsep) * \real{0.5882}}@{}}
\toprule\noalign{}
\begin{minipage}[b]{\linewidth}\raggedright
Claim
\end{minipage} & \begin{minipage}[b]{\linewidth}\raggedright
Evidence
\end{minipage} \\
\midrule\noalign{}
\endhead
\bottomrule\noalign{}
\endlastfoot
Knowledge is durable & A1 + A3 + R3: survives crashes, restarts,
conflicts \\
Cost is proportional to work & A2 + A5 + R2: no idle waste,
auto-cleanup \\
Efficiency increases & A4 + R1: \(C_{learn} \searrow\), \(C_{boot}\)
bounded \\
Multi-agent is sub-linear & A6 + R3: \(C(N) < N \times C_{solo}\),
shared project memory \\
\end{longtable}

\subsubsection{Limitations}\label{limitations}

\begin{itemize}
\tightlist
\item
  Convergence (A4) is empirical, not proven --- LLM noise can prevent it
\item
  Budget pruning may discard relevant low-priority knowledge under
  extreme pressure
\end{itemize}

\subsubsection{Future Work}\label{future-work}

\begin{itemize}
\tightlist
\item
  Formal convergence proof under bounded noise assumptions
\item
  Cross-relay federation for multi-organization agent networks
\end{itemize}

\begin{center}\rule{0.5\linewidth}{0.5pt}\end{center}

\emph{The agent does not live. It spawns, works, learns, dies. The next
one knows everything.}


    % Add a bibliography block to the postdoc
    
    
    
\end{document}
